 \setlength{\baselineskip}{20pt}
\chapter{总结与展望}
\label{cha:chap5}



\section{研究总结}

我这里更关注的是问题建模方式的变化。
第二章将停车位检测从基于规则的检测流程重构为结构感知的端到端建模问题；
第三章则在此基础上，将建模维度由单帧空间结构扩展到多帧时序一致性建模，因此并不是简单的模块叠加，而是检测范式在空间和时间维度上的逐步扩展。

创新点一：构建大规模复杂场景停车位数据集并提出半监督检测框架

针对现有停车位检测研究中高质量标注数据匮乏、复杂场景覆盖不足的问题，本文构建了一个覆盖多光照、多视角、多停车形式的大规模停车位数据集，显著提升了数据的多样性与复杂性。在此基础上，提出了一种面向停车位检测任务的半监督检测框架，通过教师–学生学习范式有效利用未标注数据，缓解了标注成本高昂对模型训练与泛化能力的制约，为后续检测方法的研究提供了统一且可靠的数据与训练基础。

创新定性关键词：
数据层创新 + 学习范式创新


创新点二：提出单帧结构感知的端到端停车位检测方法

针对传统停车位检测方法中依赖启发式规则进行后处理匹配、结构一致性难以保证的问题，本文提出了一种单帧结构感知的端到端停车位检测方法。该方法从建模角度将停车位检测问题重构为基于关键点的结构化预测问题，在端到端训练过程中显式引入停车位内部结构关系约束，使结构一致性直接参与模型优化，从而避免了规则依赖并提升了复杂与密集场景下的检测稳定性与一致性。

创新定性关键词：
问题建模重构 + 结构感知 + 端到端优化

创新点三：提出多帧时序特征融合的停车位检测方法

针对真实驾驶场景中停车位检测结果在连续帧条件下易受遮挡、视角变化和噪声影响而不稳定的问题，本文在单帧结构感知检测方法的基础上，进一步提出了一种多帧时序特征融合的停车位检测方法。该方法通过对连续帧中的结构感知表示进行时序建模，将停车位检测从静态空间结构建模扩展至时空联合建模，有效提升了检测结果在时间维度上的一致性与鲁棒性，使所提出的方法更加贴近实际应用需求。

创新定性关键词：
建模维度扩展 + 时序一致性建模


本文围绕复杂场景下停车位检测任务，从数据与学习范式、单帧结构建模以及多帧时序建模三个层面展开研究，逐步将停车位检测从基于规则的静态检测流程，发展为结构感知且具备时序一致性的端到端检测方法


\section{未来展望}
